\documentclass[10pt,a4paper]{article}

\usepackage[T1]{fontenc}
\usepackage[utf8]{inputenc}
\usepackage[portuguese]{babel}
\usepackage[margin=2.1cm,top=2.3cm,bottom=2.2cm]{geometry}
\usepackage{newtxtext,newtxmath}
\usepackage{inconsolata}
\usepackage{amsmath}
\usepackage{xcolor}
\usepackage{tcolorbox}
\usepackage{titlesec}
\usepackage{enumitem}
\usepackage{fancyhdr}
\usepackage[colorlinks=true,linkcolor=azul,urlcolor=azul]{hyperref}

\definecolor{azul}{HTML}{2A4B8D}
\definecolor{ferrugem}{HTML}{A63A26}
\definecolor{fundo}{HTML}{F4F4F0}
\definecolor{borda}{HTML}{D8D8D1}
\definecolor{cinza}{HTML}{5A5F66}
\definecolor{verde}{HTML}{2E7D4F}

\tcbuselibrary{skins,breakable}

% ---- caixa de equacao ------------------------------------------------
\newenvironment{eqn}[2]{%
  \def\eqnnum{#1}\def\eqnsrc{#2}%
  \begin{tcolorbox}[enhanced, breakable, sharp corners,
    colback=fundo, colframe=borda, boxrule=0.4pt,
    borderline west={1.6pt}{0pt}{azul},
    left=8pt, right=8pt, top=5pt, bottom=3pt,
    before skip=8pt, after skip=3pt]%
}{%
  \par\vspace{3pt}{\color{borda}\rule{\linewidth}{0.3pt}}\par\vspace{1pt}%
  {\footnotesize\bfseries\color{azul}(\eqnnum)}\hfill
  {\footnotesize\color{cinza}\ttfamily\eqnsrc}%
  \end{tcolorbox}}

\newcommand{\nota}[1]{\par\vspace{2pt}{\footnotesize\color{cinza}#1}\par}

% ---- estado de um problema -------------------------------------------
\newcommand{\chip}[2]{%
  \tcbox[on line, sharp corners, boxrule=0.5pt, colframe=#1, colback=white,
         left=3pt,right=3pt,top=1pt,bottom=1pt, box align=base]
        {\scriptsize\bfseries\color{#1}\MakeUppercase{#2}}}
\newcommand{\bloqueia}{\chip{ferrugem}{bloqueia}}
\newcommand{\divida}{\chip{azul}{dívida}}
\newcommand{\limite}{\chip{cinza}{limite}}
\newcommand{\fechado}{\chip{verde}{fechado}}

\newcommand{\prob}[3]{%
  \par\vspace{11pt}{\color{borda}\rule{\linewidth}{0.4pt}}\par\vspace{5pt}%
  \noindent\makebox[2.4em][l]{\bfseries\color{azul}#1}%
  \begin{minipage}[t]{\dimexpr\linewidth-2.4em\relax}%
  {\bfseries #2}\hspace{6pt}#3\par\vspace{3pt}}
\newcommand{\fimprob}{\end{minipage}\par}

\newcommand{\loc}[1]{\par\vspace{3pt}{\footnotesize\color{cinza}\ttfamily #1}\par}
\newcommand{\evid}[1]{%
  \par\vspace{4pt}
  \begin{tcolorbox}[enhanced, sharp corners, colback=fundo, colframe=borda,
    boxrule=0.4pt, borderline west={1.6pt}{0pt}{cinza},
    left=6pt,right=6pt,top=3pt,bottom=3pt, before skip=2pt, after skip=4pt]
  {\ttfamily\footnotesize #1}
  \end{tcolorbox}}

\titleformat{\section}{\normalfont\large\bfseries\color{azul}}{\thesection.}{0.6em}{}
\titleformat{\subsection}{\normalfont\normalsize\bfseries}{}{0em}{}
\titlespacing{\section}{0pt}{18pt}{6pt}
\titlespacing{\subsection}{0pt}{14pt}{4pt}

\pagestyle{fancy}\fancyhf{}
\renewcommand{\headrulewidth}{0.4pt}
\fancyhead[L]{\footnotesize\color{cinza}PHASIS --- equações e pendências}
\fancyhead[R]{\footnotesize\color{cinza}\thepage}

\setlength{\parindent}{0pt}
\setlength{\parskip}{4pt}
\emergencystretch=2.5em
\hbadness=10000

\newcommand{\rt}{r_{\rm t}}
\newcommand{\rs}{r_{\rm s}}
\newcommand{\Einf}{E_\infty}
\newcommand{\Eloc}{E_{\rm loc}}

\begin{document}

\begin{center}
{\footnotesize\color{cinza}\MakeUppercase{Relatório de estado}\quad$\cdot$\quad
22 de agosto de 2026\quad$\cdot$\quad revisto após a corrente neutra}\\[6pt]
{\LARGE\bfseries Equações e Pendências do PHASIS}\\[5pt]
\begin{minipage}{0.86\linewidth}\centering
Onde vive cada equação física, do dipolo de cor à geodésica nula, e os
problemas abertos --- ordenados por aquilo que impede a varredura de
produção.
\end{minipage}\\[8pt]
{\small\color{cinza}
\textbf{30} equações \quad$\cdot$\quad
\textbf{247} verificações \quad$\cdot$\quad
\textbf{0} falhas \quad$\cdot$\quad
\textbf{12} problemas abertos, 2 fechados}
\end{center}

\vspace{6pt}

% =====================================================================
\section{As equações, e onde elas estão}

Duas cadeias independentes, acopladas \emph{por arquivo} --- a tabela de
seção de choque --- e não por link. O \texttt{dipole/} produz
$\sigma(E)$; o PHASIS a consome ao longo de geodésicas.

\subsection{A. Seção de choque $\nu N$}

\[
\psi_T,\psi_L \;\longrightarrow\; \sigma_{\rm dip}
\;\longrightarrow\; F_T,F_L
\;\longrightarrow\; \frac{d^2\sigma}{dx\,dy}
\;\longrightarrow\; \sigma_{CC}(E)
\]

Formalismo de Kutak \& Kwieciński, EPJ C29 (2003) 521, eqs.~(2), (8),
(9), (12), (13).

\begin{eqn}{1}{dipole/src/wavefunctions.cpp:11 · epsilon2}
\[ \varepsilon^2(z,Q^2) = z(1-z)Q^2 + z\,m^2 + (1-z)\,\mu^2 \]
\nota{Massa transversa efetiva do dipolo quark--antiquark. $m$ e $\mu$
são as massas dos dois sabores do canal; $z$ é a fração de momento
longitudinal.}
\end{eqn}

\begin{eqn}{2}{dipole/src/wavefunctions.cpp:18 · psiT2\_pre}
\[ |\psi_T|^2 = \frac{6\alpha_{\rm EW}}{(2\pi)^2}
\Big[ (g_V^2+g_A^2)(z^2+\bar z^2)\,\varepsilon^2 K_1^2(\varepsilon r)
+ \big( g_V^2(zm+\bar z\mu)^2 + g_A^2(zm-\bar z\mu)^2 \big) K_0^2(\varepsilon r) \Big] \]
\nota{Com $\bar z = 1-z$. As Bessel modificadas vêm de
\texttt{std::cyl\_bessel\_k} (C++17); confere com Boost a
$2{,}7\times10^{-16}$, o que tirou o Boost das dependências.}
\end{eqn}

\begin{eqn}{3}{dipole/src/wavefunctions.cpp:47 · psiL2\_pre}
\begin{align*}
|\psi_L|^2 &= \frac{6\alpha_{\rm EW}}{(2\pi)^2 Q^2}
\Big[ A\,\varepsilon^2 K_1^2 + \big( g_V^2 B_V^2 + g_A^2 B_A^2 \big) K_0^2 \Big] \\
A &= g_V^2(m-\mu)^2 + g_A^2(m+\mu)^2 \\
B_{V,A} &= 2Q^2 z\bar z + (m\mp\mu)(zm \mp \bar z\mu)
\end{align*}
\nota{Não se assume simetria $z \leftrightarrow 1-z$: ela só vale quando
$m = \mu$, e falha no canal $(c,s)$.}
\end{eqn}

\begin{eqn}{4}{dipole/include/parameters.hpp · makeCouplings}
\[ \text{CC:}\quad g_V = -1,\qquad g_A = +1,\qquad
\alpha_{\rm EW} = \frac{M_W^2 G_F}{4\pi\sqrt2} \]
\nota{Canais favorecidos por Cabibbo: $u\bar d$ e $c\bar s$. Os dois
somados dão o $\Sigma(\text{carga}^2)$ efetivo $=4$ das eqs.~(12)--(13)
de KK.}
\end{eqn}

\begin{eqn}{5}{dipole/src/dipole\_models.cpp:17 · sigmaDipoleGBW}
\[ \sigma_{\rm dip}(r,x) = \sigma_0\left[1 - e^{-r^2/4R_0^2(x)}\right],
\qquad R_0^2(x) = Q_0^{-2}\left(\frac{x}{x_0}\right)^{\!\lambda} \]
\[ \sigma_0 = 29{,}12~\text{mb},\quad \lambda = 0{,}277,\quad x_0 = 4{,}1\times10^{-5} \]
\nota{Ajuste de quatro sabores de Golec-Biernat \& Wüsthoff, PRD 59
(1999) 014017. As massas efetivas (0,14 e 1,5~GeV) fazem parte do
ajuste --- não são convenção livre.}
\end{eqn}

\begin{eqn}{6}{dipole/src/dipole\_models.cpp:39 · Qs\_bCGC, amplitude\_bCGC, sigmaDipoleIIM}
\begin{align*}
Q_s^2(x,b) &= \left(\frac{x_0}{x}\right)^{\!\lambda}
\left[e^{-b^2/2B_{\rm CGC}}\right]^{1/\gamma_s},
\qquad \tau = r\,Q_s \\[3pt]
N &= N_0\left(\frac{\tau}{2}\right)^{2\gamma_{\rm eff}},\quad
\gamma_{\rm eff} = \gamma_s + \frac{\ln(2/\tau)}{\kappa\lambda Y},
\qquad \tau \leq 2 \\[3pt]
N &= 1 - e^{-A\ln^2(B\tau)}, \qquad \tau > 2 \\[3pt]
\sigma_{\rm dip} &= 2\!\int\! d^2b\,N = 4\pi\!\int\! b\,db\,N
\end{align*}
\nota{$Y = \ln(1/x)$, $\kappa = 9{,}9$ (valor LO do BFKL); $A$ e $B$
fixados por continuidade em $\tau = 2$. Rezaeian \& Schmidt, PRD 88
(2013) 074016, Tabela~II. \textbf{Aqui morava um erro de $10^6$}:
$x_0 = 0{,}00069$, e $6{,}46\times10^{-6}$ é a \emph{incerteza}, não um
multiplicador.}
\end{eqn}

\begin{eqn}{7}{dipole/src/integrals.cpp:94 · structureFunctionsTL}
\[ F_{T,L}(x,Q^2) = \frac{Q^2}{4\pi^2}\int\! d^2r \int_0^1\! dz\,
|\psi_{T,L}(r,z,Q^2)|^2\,\sigma_{\rm dip}(r,x) \]
\nota{Quadratura: $r$ em $\ln r$ sobre $[10^{-8},10^{3}]$~GeV$^{-1}$; $z$
em log pelas \emph{duas} pontas a partir de $z_{\min}=10^{-11}$. Grade
uniforme perdia o pico em $r\sim1/\varepsilon$ --- 0,025~GeV$^{-1}$ em
$Q^2 = M_W^2$, contra um passo de 2.}
\end{eqn}

\begin{eqn}{8}{dipole/src/sigma\_nuN\_core.cpp:39 · d2sigma\_core}
\[ \frac{d^2\sigma}{dx\,dy} = \frac{G_F^2 M_N E}{\pi}
\left(\frac{M_V^2}{Q^2+M_V^2}\right)^{\!2}
\left[ \frac{1+(1-y)^2}{2}F_2 - \frac{y^2}{2}F_L
\pm y\!\left(1-\frac{y}{2}\right)xF_3 \right] \]
\nota{Sinal $+$ para $\nu$, $-$ para $\bar\nu$.
$F_2 = (1-x)^7(F_T+F_L)$ --- regra de contagem de constituintes com
$n_s = 4$ espectadores. $xF_3$ vem do LHAPDF, não do dipolo.
$M_V = M_W$ (CC) ou $M_Z$ (NC): é a única constante que separa as duas
correntes.}
\end{eqn}

\begin{eqn}{9}{dipole/src/sigma\_nuN\_core.cpp:86 · sigmaNuN}
\[ \sigma(E) = \int_{Q^2_{\min}}^{0{,}999\,s}\!\!dQ^2
\int_{Q^2/s}^{1}\!dx\;\frac{1}{x\,s}\,\frac{d^2\sigma}{dx\,dy},
\qquad y = \frac{Q^2}{x s},\quad s = 2M_N E \]
\nota{Simpson composto em $\ln Q^2$ e $\ln x$, com \emph{densidade} de
nós por década fixa (não número fixo), e $N_{\ln x}$ constante ao longo
do laço de $Q^2$. \textbf{O extremo inferior em $x$ é exatamente
$y = 1$} --- ver a nota de fechamento da Parte II.}
\end{eqn}

A tabela \textbf{colinear} usa exatamente (8) e (9), com $F_2$, $F_L$ e
$xF_3$ do yadism NLO sobre NNPDF3.1, e \emph{sem} o fator $(1-x)^7$ ---
os PDFs colineares já se anulam em $x\to1$ sozinhos.
\texttt{tools/make\_collinear\_table.py}.

\subsection{B. Transporte ao longo da geodésica}

\begin{eqn}{10}{include/phasis/metric.hpp:18}
\[ ds^2 = -f(r)\,dt^2 + h(r)\,dr^2 + r^2 d\Omega^2,
\qquad \text{Schwarzschild: } f = 1-\frac{\rs}{r},\; h = \frac{1}{f} \]
\nota{A interface expõe só $f$ e $h$. Nada a jusante supõe $hf = 1$ ---
em Schwarzschild o produto some, em métricas gerais não.}
\end{eqn}

\begin{eqn}{11}{src/metric.cpp:8 · Metric::E\_local}
\[ \Eloc(r) = \frac{\Einf}{\sqrt{f(r)}} \]
\nota{A energia pedida a $\sigma$ é sempre a \emph{local}, medida pelo
observador estático no ponto. Fica \emph{dentro} da integral de $\tau$,
mesmo na Fase~1 onde $f=1$ e fatoraria --- é o que o teste T10
verifica.}
\end{eqn}

\begin{eqn}{12}{src/metric.cpp:64 · Metric::r\_turning ; Schwarzschild::r\_turning}
\[ g(r) \equiv \frac{r}{\sqrt{f(r)}} = b
\qquad\Longrightarrow\qquad
\rt = \frac{2b}{\sqrt3}\cos\!\left[\frac{1}{3}\arccos\!\left(-\frac{3\sqrt3\,\rs}{2b}\right)\right] \]
\nota{A rota genérica bissecta $g(r)-b$, descendo a partir de
$r_{\max}$ e parando quando $f$ deixa de ser positivo. As duas rotas se
validam uma contra a outra.}
\end{eqn}

\begin{eqn}{13}{src/metric.cpp · Schwarzschild::is\_captured, r\_photon\_sphere}
\[ b_{\rm crit} = \frac{3\sqrt3}{2}\rs = g(1{,}5\,\rs),
\qquad \text{captura} \iff 4b^2 \leq 27\rs^2,
\qquad r_{\rm ph} = 1{,}5\,\rs \]
\nota{$g$ tem \emph{mínimo} na esfera de fótons. É por isso que a
topologia do raio emitido tem quatro casos e não três: abaixo de
$r_{\rm ph}$ a condição de escape inverte.}
\end{eqn}

\begin{eqn}{14--15}{src/metric.cpp:17 · turning\_factor, turning\_factor\_at\_turn}
\begin{align*}
w(r) &= 1 - \frac{f(r)\,b^2}{r^2} \qquad\text{(anula-se em $\rt$)}\\[3pt]
T(r) &= \frac{f(\rt)\,r^2 - f(r)\,\rt^2}{r-\rt},
\qquad T(\rt) = 2f(\rt)\rt - f'(\rt)\rt^2 \\[3pt]
T_{\rm Schw} &= (r+\rt) - \frac{\rs}{\rt r}\left(r^2 + r\rt + \rt^2\right)
\end{align*}
\nota{A fatoração central do projeto:
$w = (r-\rt)\,T(r)\,/\,\big(f(\rt)r^2\big)$, com $T$ regular e positiva
em $\rt$.}
\end{eqn}

\begin{eqn}{16}{src/trace.cpp:121 · Geometry::dl\_ds ; ray\_derivatives}
\[ \text{com } r = \rt + s^2:\qquad
\frac{dl}{ds} = 2r\sqrt{\frac{h(r)\,f(\rt)}{T(r)}} \]
\nota{O fator $s$ do jacobiano cancela \emph{analiticamente} a raiz de
$w$ --- a singularidade integrável do ponto de retorno some antes de
qualquer aritmética de ponto flutuante. Sobra a divergência
\emph{logarítmica} quando $b\to b_{\rm crit}^+$, que é física: o raio
enrola na esfera de fótons.}
\end{eqn}

\begin{eqn}{17}{src/trace.cpp:193 · half\_sweep\_angle ; ray\_derivatives}
\[ \frac{d\varphi}{dr} = \frac{b}{r^2}\frac{\sqrt{h f}}{\sqrt w},
\qquad\qquad \frac{d\psi}{dl} = \frac{b\sqrt f}{r^2} \]
\nota{O $\sqrt f$ não pode faltar: como $1/\sqrt f \approx 1+\rs/2r$,
omiti-lo injeta um termo da \emph{mesma ordem} da resposta e dá
$\Delta\varphi = 3\rs/b$ em vez de $2\rs/b$. Foi um bug real; T8 depois
pegou o coeficiente de segunda ordem, $15\pi/16$.}
\end{eqn}

\begin{eqn}{18}{src/trace.cpp:161 · Integrand::operator() ; trace\_ray}
\[ \tau = \int N_A\,\rho(r,\theta)\,\sigma\big(\Eloc(r)\big)\,dl,
\qquad P_{\rm surv} = e^{-\tau},
\qquad \cos\theta = \sin i \,\sin\psi \]
\nota{Movimento planar: basta girar o plano orbital. $i$ é a inclinação
em relação ao equador; $\psi$ a posição azimutal dentro do plano, medida
da linha de nodos. Perfil esférico vai por quadratura adaptativa;
perfil com $\theta$ exige a rota de EDO acoplada --- ver \textbf{P1}.}
\end{eqn}

\begin{eqn}{19}{src/trace.cpp:105 · Geometry::Mode}
\begin{align*}
\text{NoTurn (emitido para fora, } b<b_{\rm crit}):&\quad
\frac{dl}{ds} = 2s\sqrt{\frac{h}{w}},\quad \text{âncora } r_{\rm emit}\\
\text{Radial } (b=0):&\quad \frac{dl}{ds} = 2s\sqrt{h}
\end{align*}
\nota{Nesse regime o ponto de retorno \emph{não existe} --- não é
virtual, é ausente, porque $b < \min g$. Mas $w$ nunca chega perto de
zero, então não há singularidade e não há o que fatorar. O caso radial
precisa de rota própria porque $\rt = 0$ e $f(0)$ são ambos
singulares.}
\end{eqn}

\begin{eqn}{20}{src/emission.cpp:12 · b\_from\_emission}
\[ b = \frac{r_{\rm emit}\sin\psi}{\sqrt{f(r_{\rm emit})}} \]
\nota{No referencial ortonormal estático, com $\psi$ medido da direção
radial para fora: $L = r_{\rm emit}\Eloc\sin\psi$ e
$\Einf = \sqrt f\,\Eloc$. Auto-verificação: em $\psi = \pi/2$ dá
$b = g(r)$, ou seja $\rt = r_{\rm emit}$. Emissão isotrópica local é
\textbf{$\cos\psi$ uniforme}, não $\psi$ uniforme.}
\end{eqn}

\begin{eqn}{21}{src/emission.cpp:52 · escape\_cone\_angle, captured\_fraction}
\[ \sin\psi_c(r) = \frac{b_{\rm crit}\sqrt{f(r)}}{r},
\qquad f_{\rm cap} = \frac{1-\cos\psi_c}{2} \]
\nota{Dois valores exatos, verificados a 0 e $1{,}4\times10^{-16}$:
$r = 1{,}5\rs \Rightarrow \psi_c = \pi/2$ e $f_{\rm cap} = 1/2$;
$r = 3\rs \Rightarrow \psi_c = \pi/4$ e $f_{\rm cap} = (2-\sqrt2)/4$.}
\end{eqn}

\begin{eqn}{22}{src/cascade.cpp:296 · CascadeCtx ; CascadeKernel}
\begin{align*}
\frac{d\phi_i}{dl} &= -n(l)\,\sigma_{\rm tot}(\Eloc^{(i)})\,\phi_i
+ n(l)\sum_{j\geq i}\sigma_{NC}(\Eloc^{(j)})\,G_{ij}\,\phi_j \\[3pt]
&\qquad \sum_i G_{ij} + G_{{\rm leak},j} = 1
\end{align*}
\nota{$y$ é \textbf{invariante sob redshift}: o neutrino inicial e o
final estão no mesmo $r$, o mesmo $\sqrt f$ aparece nos dois, logo
$\Eloc/\Eloc' = \Einf/\Einf' = 1-y$. A \emph{estrutura} $G_{ij}$ é
montada uma vez e nunca mais; só o coeficiente varia ao longo do raio. O
vazamento é rastreado numa componente própria, nunca empilhado no bin de
baixo.}
\end{eqn}

\begin{eqn}{23}{src/cross\_section.cpp:52 · log\_slope\_detail}
\[ \alpha_{\rm eff}(E) = \frac{d\ln\sigma}{d\ln E},
\qquad \text{teto}(E) = 3^{\alpha_{\rm eff}/2} \]
\nota{O teto vem de $\Eloc/\Einf = 1/\sqrt{f(\rt)} \leq \sqrt3$: todo
raio não capturado tem $\rt \geq 1{,}5\rs$, onde $f = 1/3$. A assinatura
observável não é o valor do teto --- é a \emph{derivada} dele em $E$.}
\end{eqn}

\begin{eqn}{24}{src/redshift\_scan.cpp ; trace\_ray(..., freeze\_redshift)}
\begin{align*}
R(b,E) &= \frac{\tau_{\rm full} - \tau_{\rm frozen}}{\tau_{\rm frozen}}
= \big\langle f^{-\alpha/2}\big\rangle - 1,
\qquad \text{peso } n(r)\sigma\,dl \\[3pt]
\text{teto} - 1 - R &= \frac{K}{\tau_{\rm frozen}},
\qquad K \text{ constante}
\end{align*}
\nota{Identidade, não ajuste: $\int w\,(g_{\max}-g)$ converge porque
$(g_{\max}-g)$ se anula exatamente onde o peso diverge, enquanto
$\int w$ diverge logaritmicamente. Dois pontos extrapolam
$3^{\alpha/2}$ a $3\times10^{-11}$. $R$ é invariante sob
$\sigma \to k\sigma$ --- é isso que legitima comparar dipolo com
colinear apesar da supressão global do primeiro.}
\end{eqn}

\subsection{C. Corrente neutra}

Camada mais nova. O CC junta dois sabores num dipolo; o $Z$ acopla ao
\emph{mesmo} sabor dos dois lados, e é daí que vem a maior parte da
razão $\sigma_{NC}/\sigma_{CC}$ --- não do acoplamento.

\begin{eqn}{25}{dipole/include/parameters.hpp · makeCouplings}
\[ g_V^q = T_3 - 2Q_q\sin^2\theta_W,
\qquad g_A^q = T_3,
\qquad C = \frac{g_Z}{2} = \frac{g}{2\cos\theta_W} \]
\[ \alpha_{NC} = \frac{C^2}{4\pi} = \frac{\sqrt2\,G_F M_Z^2}{4\pi} \]
\nota{\textbf{D10, resolvido.} O valor antigo era metade deste --- mas
\emph{não muda} $\sigma$, porque a tabela divide $F_T$ e $F_L$ pelo
$\alpha_{\rm EW}$ do próprio canal (convenção DIS). As duas metades do
comentário antigo eram fatores 2 diferentes, e ambas verdadeiras. A
mesma rota dá $C_{CC} = g/(2\sqrt2)$ e reproduz o $\alpha_{CC}$ que já
estava certo.}
\end{eqn}

\begin{eqn}{26}{dipole/src/sigma\_nuN.cpp · channels\_nc}
\begin{align*}
\text{CC:}&\quad (u\bar d),\,(c\bar s) \quad\text{--- 2 dipolos} \\
\text{NC:}&\quad (u\bar u),\,(d\bar d),\,(s\bar s),\,(c\bar c)
\quad\text{--- 4 dipolos} \\[4pt]
&\Sigma_{NC}(g_V^2+g_A^2) = 1{,}3127,
\qquad \Sigma_{CC} = 2\times2 = 4 \\[3pt]
&\frac{\Sigma_{NC}}{\Sigma_{CC}}\left(\frac{M_Z}{M_W}\right)^{\!2} = 0{,}4224
\end{align*}
\nota{$1{,}3127$ é o $C$ de KK eq.~(15). Mesma população de quarks nos
dois casos ($u,d,s,c$ --- o ajuste GBW é de quatro sabores), contagem
diferente de dipolos. Medido: 0,305 em $10^3$~GeV subindo a 0,412 em
$10^{13}$, com o assintótico 3\% abaixo de 0,4224 por efeito de massa do
charme.}
\end{eqn}

\begin{eqn}{27}{dipole/src/sigma\_nuN\_core.cpp · d2sigma\_core, bosonMass2}
\[ \text{A eq.~(8) vale para as \emph{duas} correntes, com }
M_V = M_W \text{ ou } M_Z \]
\nota{Uma função só, não duas. Não é economia de linhas: a razão
$\sigma_{NC}/\sigma_{CC}$ só significa algo se o resto for idêntico, e
aqui é o mesmo código. Verificado: depois do refactor a tabela CC saiu
\textbf{bit a bit idêntica} em 300 energias.}
\end{eqn}

\begin{eqn}{28}{dipole/src/weak\_structure\_functions.cpp:63 · xF3\_NC\_isoscalar}
\[ xF_3^{NC} = 2\sum_q g_V^q g_A^q \left( xq - x\bar q \right) \]
\nota{De $g_L^2 - g_R^2 = g_V g_A$. É puramente de valência, e cai muito
mais rápido que o total: acima de $\sim10^6$~GeV contribui menos de 1\%.
O $xF_3$ do \emph{CC}, por comparação, vale $+70\%$ de $\sigma$ em
$10^3$~GeV e $+0{,}16\%$ em $10^{14}$ --- e é por isso que
$\sigma_{NC}/\sigma_{CC}$ sobe com a energia.}
\end{eqn}

\begin{eqn}{29}{dipole/src/sigma\_nuN\_core.cpp · dsigma\_dy, yMinKinematic}
\[ \frac{d\sigma}{dy} = \int_{Q^2_{\min}/ys}^{x_{\max}}\!\!dx\;
\frac{d^2\sigma}{dx\,dy},
\qquad Q^2 = x\,y\,s,
\qquad \text{suporte: } y \geq \frac{Q^2_{\min}}{s\,x_{\max}} \]
\nota{O corte inferior em $y$ \textbf{depende da energia} (cai como
$1/E$), mas o formato de tabela exige uma grade em $y$ comum. A grade usa
o corte da energia \emph{mais alta}, e nas energias baixas as linhas
abaixo do próprio corte saem zeradas --- zeros de verdade, não lacunas.
Ordenar a mesma integral dupla de dois jeitos é teste de consistência:
$\int dy\,(d\sigma/dy)$ recupera $\sigma$ em 1,04\%, que é o erro do
trapézio sobre 120 nós.}
\end{eqn}

\begin{eqn}{30}{src/cascade.cpp · CascadeKernel}
\[ K_{ij}(E) = \sigma_{NC}(E)\,G_{ij}(E),
\qquad \sum_i G_{ij}(E) + G_{{\rm leak},j}(E) = 1 \]
\[ G \text{ tabelado em nós log-}E\text{, interpolado \textbf{linearmente} em } \ln E \]
\nota{A forma em $y$ da tabela real varia por \textbf{fator 40} entre
$10^5$ e $10^{11}$~GeV, então $G_{ij}$ constante não serve. A
linearidade não é conveniência: a identidade de consistência é
\emph{linear} nos $G$, logo interpolação linear de um conjunto que soma
1 em cada nó ainda soma 1 \emph{entre} os nós. A cascata conserva número
exatamente em qualquer $E$. Medido $2{,}22\times10^{-16}$ nos nós e
entre eles. Os intervalos em $y$ não dependem de $\Eloc$ --- $y$ é
invariante sob redshift e as bordas dos bins estão em $\Einf$ --- e é
por isso que dá para tabelar uma vez e nunca reintegrar por passo.}
\end{eqn}

Fora dessa lista, e sem equação própria: \texttt{src/shell.cpp} localiza
a casca $\tau = 1$ com um protocolo que não presume monotonicidade ---
varredura grossa de $\geq32$ pontos, contagem de mudanças de sinal,
bisseção de \emph{todos} os brackets, e \texttt{shell\_found = false}
limpo quando não há nenhum.

% =====================================================================
\section{Os problemas}

Quatorze, ordenados pelo que impede o próximo passo. \bloqueia{} impede
a varredura de produção; \divida{} é defeito conhecido e documentado;
\limite{} é limitação honesta, não defeito; \fechado{} foi resolvido
depois da primeira versão deste relatório.

\prob{P1}{A rota de EDO nunca recebeu a refatoração de emissão}{\bloqueia}
A topologia do raio emitido (T31--T35) foi aplicada só à rota de
quadratura. A rota de EDO --- \emph{obrigatória} para perfil não
esférico, que é justamente o disco --- continua tratando todo raio como
vindo do infinito: integra os dois ramos sobre a faixa inteira, ignora
\texttt{r\_emit}, e usa \texttt{geo.r\_turn} onde deveria usar
\texttt{geo.r\_anchor}.

Medido com perfil esférico e \texttt{force\_ode}, contra a rota de
quadratura que está validada:
\evid{raio do infinito, b = 5 r\_s\ \ \ \ \ tau\_quad = 79,1076\ \ \ tau\_ode = 79,1076\ \ \ razao 1,0000000\\
emitido r=10 r\_s, psi=30$^\circ$\ \ \ \ \ \ \ tau\_quad = 10,2996\ \ \ tau\_ode = 73,5825\ \ \ razao 7,144\\
emitido r=1,3 r\_s, psi=10$^\circ$\ \ \ \ \ \ tau\_quad = 112,173\ \ \ tau\_ode = LANCOU excecao\\
emitido r=10 r\_s, psi=0$^\circ$\ \ \ \ \ \ \ \ tau\_quad = \ 9,8136\ \ \ tau\_ode = \ 0\ \ \ \ \ \ \ \ razao 0}
E no caso astrofísico de verdade --- \texttt{FlaredThinDisk}, que força a
rota de EDO sem nenhuma flag --- $\psi=10^\circ$ lança,
$\psi=30^\circ$ produz um \emph{ramo de entrada que não deveria existir}
(\texttt{tau\_in} = 0,0134) com $r_{\min} = 4{,}67\rs$ \textbf{abaixo} do
raio de emissão de $10\rs$, e $\psi=0^\circ$ dá zero em silêncio.

O caso radial é o pior dos três: erra \emph{em silêncio}. Nenhum teste
pegou porque T31--T35 usaram só perfis esféricos.
\loc{src/trace.cpp:304, :558 · trace\_via\_ode, OdeCtx}
\fimprob

\prob{P2}{A cascata NC herda o mesmo defeito}{\bloqueia}
\texttt{transport\_cascade} não menciona \texttt{r\_emit},
\texttt{outward} nem \texttt{classify} em lugar nenhum, e o seu
\texttt{CascadeCtx} também calcula \texttt{r = r\_turn + s*s}.
Transportar um espectro a partir de um emissor a distância finita hoje
dá o mesmo tipo de resposta errada de P1.

Ficou mais concreto desde a primeira versão deste relatório: a cascata
agora roda com tabela real, então o defeito dela deixou de ser
hipotético.

Os dois têm a mesma correção estrutural: a escolha de âncora e o recorte
por ramo precisam sair de \texttt{trace\_ray} e virar propriedade
compartilhada da \texttt{Geometry}, consumida pelas três rotas.
\loc{src/cascade.cpp:296 · CascadeCtx::operator()}
\fimprob

\prob{P3}{D10 --- o acoplamento NC do dipolo está um fator 2 baixo}{\fechado}
O valor \emph{estava} errado: $\alpha_{NC} = \sqrt2 G_F M_Z^2/4\pi$, o
dobro do que havia. Mas o comentário antigo tinha duas metades que
pareciam se contradizer, e ambas eram verdadeiras --- eram
\textbf{fatores 2 diferentes}. A tabela divide $F$ pelo
$\alpha_{\rm EW}$ do próprio canal, então $\alpha_{NC}$ \textbf{cancela}
e $\sigma$ não se mexe. Por isso a checagem numérica antiga dava o valor
certo do SM apesar do valor errado da constante: ela mede a razão das
somas de carga, onde $\alpha$ já saiu.

Ver eq.~(25). Corrigido e o comentário reescrito para explicar por que
não muda $\sigma$ --- sem isso, alguém ``consertaria'' $\sigma_{NC}$ por
um fator 2 mais adiante.
\loc{dipole/include/parameters.hpp · makeCouplings, caso NC}
\fimprob

\prob{P4}{O oráculo Python divergiu do C++}{\divida}
O conserto do extremo $y=1$ entrou no C++ e não no oráculo, que ainda
tem \texttt{if not (0.0 < y < 1.0): continue}.

\textbf{É pior do que a primeira versão deste relatório dizia.} Medido:
$7{,}628\times10^{-33}$ contra $7{,}352\times10^{-33}$ em $10^9$~GeV,
3,8\% de diferença. Mas a comparação hoje não isola nada, porque o
oráculo está atrás de \emph{duas} mudanças estruturais, não uma: a regra
de densidade de nós por década --- a quadratura dele ainda usa número
fixo --- e o conserto do $y=1$. E não conhece NC.

Esse arquivo existe para provar que uma mudança no C++ fez exatamente o
que se pretendia --- foi o instrumento que fechou a campanha de correção
a $2\times10^{-6}$. Deixá-lo divergente logo \emph{depois} de ele ter
servido aposenta o instrumento.
\loc{dipole/tests/oracle/dipole\_oracle.py:174}
\fimprob

\prob{P5}{Sete arquivos de dados do dipolo estão obsoletos}{\divida}
\texttt{tau\_profile.dat}, \texttt{tau\_profile\_collapsar.dat},
\texttt{tau\_profiles\_scan\{,\_GBW,\_IIM\}.dat},
\texttt{optical\_depth\_CC\_GBW.dat}, \texttt{test\_tau\_1e12.dat} ---
todos calculados com o $\sigma$ enviesado em 0,23\%. Ou regerar, ou
marcar o viés no cabeçalho. Nenhum deles alimenta o PHASIS hoje.
\loc{dipole/data/}
\fimprob

\prob{P6}{Não existe corrente neutra no dipolo}{\fechado}
Existe agora, de ponta a ponta: funções de estrutura NC (4 dipolos de
mesmo sabor), \texttt{xF3\_NC\_isoscalar}, \texttt{dsigma\_dy}, e as
tabelas \texttt{sigma\_nuN\_NC\_*.dat} e \texttt{dsigma\_dy\_NC\_*.dat}
com as nove chaves. Eqs.~(25)--(30).

Fechá-lo abriu e fechou também a \textbf{``Fase 4b''}: o
\texttt{CascadeKernel} lançava para qualquer forma em $y$ dependente de
$E$, e a tabela real depende por um fator 40. Construído o caminho com
$G_{ij}(E)$ interpolado linearmente --- ver eq.~(30).

Ponta a ponta: sem absorção CC, $\Sigma\phi + $ vazamento conserva a
$4{,}4\times10^{-16}$; com CC, a regeneração aumenta o fluxo
sobrevivente em \textbf{$+4{,}5\%$}.
\loc{25 verificações novas em tests/test\_phase8.cpp (T44--T50)}
\fimprob

\prob{P7}{A infraestrutura da Fase 5 está pela metade}{\bloqueia}
Faltam três coisas que uma varredura de produção precisa: \textbf{T30}
(determinismo serial contra \texttt{OMP\_NUM\_THREADS} em
$\{1,2,4,8\}$, bit a bit), o \textbf{hash do conteúdo} das tabelas no
cabeçalho --- hoje vai o caminho, as nove chaves e o commit, mas não o
hash --- e \textbf{checkpointing} com escrita incremental,
\texttt{fsync} e restart.

O reboot mostrou que restart não é luxo: o gerador colinear fatiado
sobreviveu porque grava fatia por fatia. Se tivesse sido a varredura,
teria custado a corrida inteira.
\fimprob

\prob{P8}{Falta a saída CSV do espectro transmitido}{\divida}
Pendência da Fase 4. A cascata calcula o espectro na saída e o vazamento
acumulado, mas não há executável que os grave. É o produto que a frente
3 vai querer plotar contra o ângulo de visada --- e o cálculo já roda,
falta gravar.
\fimprob

\prob{P9}{A tabela colinear extrapola o PDF acima de $\sim10^{12}$~GeV}{\limite}
Abaixo de $x = 10^{-9}$ o LHAPDF continua o PDF por lei de potência. A
fração de $\sigma$ que vem dali está registrada numa quarta coluna da
tabela: \textbf{0,0\%} até $2{,}3\times10^{10}$~GeV, \textbf{0,5\%} em
$1{,}9\times10^{11}$, \textbf{4,3\%} em $1{,}6\times10^{12}$ e
\textbf{20,7\%} em $1{,}3\times10^{13}$.

Onde essa fração é grande, $\alpha_{\rm eff}$ colinear é
\emph{propriedade da extrapolação}, não medida --- o que é precisamente
a hipótese que o dipolo substitui. Não é defeito escondido, mas quem ler
o arquivo daqui a seis meses precisa saber onde ela começa a mandar.
\fimprob

\prob{P10}{Emissor orbitando --- a aberração do fluido não existe}{\limite}
Todo emissor é \textbf{estático}. A aberração entre o referencial do
fluido e o estático é o \emph{maior} dos dois efeitos:
$\beta_{\rm orb} = \sqrt{M/r}\,/\sqrt{1-2M/r}$ vale exatamente
$\mathbf{1/2}$ na ISCO ($r = 6M$), com $\gamma = 2/\sqrt3$, e o fator
Doppler $\gamma(1\pm\beta)$ varre $1/\sqrt3$ a $\sqrt3$ ---
\textbf{fator 3 em energia}, contra o teto $\sqrt3$ do redshift
gravitacional.

Adiá-lo é sequenciamento, não a afirmação de que seja pequeno.
Introduzi-lo antes significa pôr um fator 3 por cima de um efeito de 3
pontos percentuais que agora está medido.
\fimprob

\prob{P11}{A interpolação da tabela de $\sigma$ é C$^0$}{\limite}
Linear em $(\ln E, \ln\sigma)$. Reproduz uma lei de potência
\emph{exatamente}, o que é a razão da escolha, mas a derivada analítica
dela é uma função escada --- daí o estimador de $\alpha_{\rm eff}$ ser
diferença centrada, e não analítico.

Hoje não custa nada: depois do conserto do extremo $y=1$ o ruído da
tabela caiu para $7{,}7\times10^{-7}$ e o teto pontual é monótono em 0 de
89 pontos. Vira limite se alguém quiser a segunda derivada.
\loc{src/cross\_section.cpp · TableCrossSection::sigma\_tot}
\fimprob

\prob{P12}{O cenário físico da varredura não está decidido}{\bloqueia}
A topologia está fixada --- emissor estático em \texttt{r\_emit},
isotrópico no referencial local. Falta a distribuição de
\texttt{r\_emit}, o perfil do halo e o espectro injetado. Nenhuma das
outras pendências depende disso, mas o produto final depende.
\fimprob

\prob{P13}{A grade em $y$ da tabela diferencial resolve $\sigma$ a 1\%}{\limite}
O trapézio sobre os 120 nós log-espaçados recupera $1{,}0104$ de
$\sigma_{NC}$ em $10^{14}$~GeV e $0{,}9988$ em $10^3$. O 1\% é erro de
quadratura da própria tabela, não do kernel.

Não contamina a conservação: o kernel normaliza pela \textbf{soma dos
pedaços}, calculada com o mesmo integrador adaptativo que calculou cada
pedaço, e não por $\sigma_{nc}$. A discrepância entre os dois
integradores fica exposta em \texttt{worst\_norm\_mismatch()} em vez de
virar fluxo criado ou destruído ao longo do raio.

Some com mais nós em $y$, ou com trapézio em $\log y$ no leitor. Custa
1\% na normalização de $\sigma_{NC}$ dentro da cascata.
\loc{src/cascade.cpp · CascadeKernel::worst\_norm\_mismatch}
\fimprob

\prob{P14}{$xF_3$ é lixo abaixo do grid do PDF, mas pesa menos de 1\%}{\limite}
\texttt{xF3\_CC\_isoscalar} subtrai PDFs extrapolados: em
$x = 10^{-14}$ dá $\mathbf{+12{,}0}$ em $Q^2 = 1$ e $\mathbf{-10{,}6}$ em
$Q^2 = M_W^2$. O cancelamento $u - \bar u$ não sobrevive à continuação
por lei de potência abaixo de $x = 10^{-9}$.

Medido o que isso custa em $\sigma_{CC}$, que é o que importa:
\evid{E = 1e3\ \ \ +69,95\%\ \ \ \ \ \ E = 1e9\ \ \ +0,65\%\\
E = 1e5\ \ \ +11,71\%\ \ \ \ \ \ E = 1e12\ \ +0,28\%\\
E = 1e6\ \ \ \ +3,77\%\ \ \ \ \ \ E = 1e14\ \ +0,16\%}
Onde o $xF_3$ pesa (baixa energia), o $x$ típico é grande e o PDF é
medido. Onde o PDF é extrapolado (alta energia), o $xF_3$ já não pesa.
Os dois regimes não se cruzam, e o resíduo fica abaixo de 0,7\% acima de
$10^9$~GeV --- mas o número está aqui porque ninguém deveria descobrir
isso sozinho.

De quebra: o comentário em \texttt{weak\_structure\_functions.hpp} está
desatualizado --- diz 0,99\% em $10^{14}$~GeV, medido com CT10nlo antes
do conserto do extremo $y=1$. O valor agora é 0,16\%.
\loc{dipole/src/weak\_structure\_functions.cpp:36}
\fimprob

\vspace{10pt}
{\color{borda}\rule{\linewidth}{0.4pt}}

\vspace{6pt}
{\small\color{cinza}
\textbf{Resolvido, e registrado aqui porque a classe é recorrente: o
extremo $y = 1$.} O laço em $x$ começa em $x = Q^2/s$, que é exatamente
$y = 1$ --- extremo cinemático com valor \emph{finito}, não ponto
proibido. Zerá-lo punha um degrau no nó de borda e fazia Simpson errar
em $O(h)$ em vez de $O(h^4)$. Ruído de $\ln\sigma$ caiu de
$8{,}2\times10^{-4}$ para $1{,}2\times10^{-6}$, e o viés de 0,23\% em
$\sigma$ sumiu. É a \textbf{terceira ocorrência} da mesma classe:
fronteira dura avaliada exatamente num extremo de integração --- as duas
anteriores foram bordas de suporte de densidade, nas Fases 3 e 4.
\par\vspace{6pt}
\textbf{E um erro meu, registrado porque o que o pegou importa.} A
primeira tabela NC deu $\sigma_{NC}/\sigma_{CC} = \mathbf{2{,}49}$ em vez
de 0,42. Os canais eram \texttt{uu, dd, ss, cc} mas os acoplamentos
continuaram CC, porque o argumento \texttt{current} não foi passado ao
construtor da tabela --- e ele tinha \textbf{valor padrão}. O resultado
era liso, monotônico, sem exceção, com $\alpha_{\rm eff}$
bem-comportado. Nada interno ao código denunciou; quem denunciou foi o
\textbf{teste de física}. O conserto estrutural foi tirar o default, e o
compilador apontou na hora o segundo call site que eu não teria lembrado
de olhar. Um default silencioso num parâmetro que muda a física é um
convite.}

% =====================================================================
\section{O que desbloqueia o quê}

A corrente neutra saiu da frente: P3 e P6 estão fechados, e com eles a
``Fase 4b'' que o kernel deferia. O que bloqueia a varredura de produção
não mudou.

\begin{enumerate}[leftmargin=*,itemsep=3.5pt,topsep=3pt]
\item \textbf{P1 + P2} continuam sendo o bloqueio real. São o mesmo
defeito estrutural --- a rota de EDO e a cascata nunca receberam a
topologia de emissão --- e devem ser consertados juntos, com um teste que
confronte as três rotas no mesmo raio emitido. A ausência desse teste é o
que deixou passar.

\item \textbf{P7} entra dentro da frente 3, não antes: uma varredura de
produção sem hash de tabela e sem restart é uma varredura em que não se
pode confiar nem retomar.

\item \textbf{P4 e P5} continuam curtos e continuam fechando a campanha
do dipolo. P4 ficou mais urgente: o oráculo está atrás de \emph{dois}
consertos, e era ele que validava mudanças como as de hoje.

\item \textbf{P8} --- a saída CSV do espectro transmitido --- é o que
falta para a cascata NC virar produto.

\item \textbf{P13 e P14} são limites medidos, não defeitos. Ficam
registrados para que ninguém os redescubra do zero.

\item \textbf{P10} por último, como sistemática, pelo argumento da
própria magnitude.
\end{enumerate}

\vspace{10pt}{\color{borda}\rule{\linewidth}{0.4pt}}\par\vspace{4pt}
{\footnotesize\color{cinza}
PHASIS $\cdot$ commit \texttt{4314c3e} $\cdot$ 247 verificações em 8
fases, 0 falhas $\cdot$ dipolo e transporte acoplados por arquivo, não
por link. Os comandos para reproduzir cada número estão em
\texttt{docs/comandos.pdf}.}

\end{document}
